\title{2021 年高考数学 新高考I卷}
\date{}
\maketitle

\section*{一、单项选择题}
本题共8小题，每小题5分，共40分。在每小题给出的四个选项中，只有一项是符合题目要求的。

\begin{question}[points=5]
设集合 $A = \{x \mid -2 < x < 4\}$, $B = \{2,3,4,5\}$, 则 $A \cap B =$ ( )
\begin{choices}
  \item $\{2\}$
  \item $\{2,3\}$
  \item $\{3,4\}$
  \item $\{2,3,4\}$
\end{choices}
\begin{solution}
\textbf{答案：}$\{2,3\}$

\textbf{分析：}利用交集的定义可求 $A \cap B$.

\textbf{详解：}由题设有 $A \cap B = \{2,3\}$, 故选B.
\end{solution}
\end{question}

\begin{question}[points=5]
已知 $z = 2 - i$, 则 $z(\overline{z} + i) =$ ( )
\begin{choices}
  \item $6 - 2i$
  \item $4 - 2i$
  \item $6 + 2i$
  \item $4 + 2i$
\end{choices}
\begin{solution}
\textbf{答案：}$6 + 2i$

\textbf{分析：}利用复数的乘法和共轭复数的定义可求得结果.

\textbf{详解：}因为 $z = 2 - i$, 故 $\overline{z} = 2 + i$, 故 $z(\overline{z} + i) = (2 - i)(2 + 2i) = 6 + 2i$, 故选C.
\end{solution}
\end{question}

\begin{question}[points=5]
已知圆锥的底面半径为 $\sqrt{2}$, 其侧面展开图为一个半圆, 则该圆锥的母线长为 ( )
\begin{choices}
  \item $2$
  \item $2\sqrt{2}$
  \item $4$
  \item $4\sqrt{2}$
\end{choices}
\begin{solution}
\textbf{答案：}$2\sqrt{2}$

\textbf{分析：}设圆锥的母线长为 $l$, 根据圆锥底面圆的周长等于扇形的弧长可求得 $l$ 的值.

\textbf{详解：}设圆锥的母线长为 $l$, 由于圆锥底面圆的周长等于扇形的弧长, 则 $\pi l = 2\pi \times \sqrt{2}$, 解得 $l = 2\sqrt{2}$, 故选B.
\end{solution}
\end{question}

\begin{question}[points=5]
下列区间中, 函数 $f(x) = 7\sin\left(x - \frac{\pi}{6}\right)$ 单调递增的区间是 ( )
\begin{choices}
  \item $\left(0, \frac{\pi}{2}\right)$
  \item $\left(\frac{\pi}{2}, \pi\right)$
  \item $\left(\pi, \frac{3\pi}{2}\right)$
  \item $\left(\frac{3\pi}{2}, 2\pi\right)$
\end{choices}
\begin{solution}
\textbf{答案：}$\left(0, \frac{\pi}{2}\right)$

\textbf{分析：}解不等式 $2k\pi - \frac{\pi}{2} < x - \frac{\pi}{6} < 2k\pi + \frac{\pi}{2} (k \in \mathbb{Z})$, 利用赋值法可得出结论.

\textbf{详解：}因为函数 $y = \sin x$ 的单调递增区间为 $\left(2k\pi - \frac{\pi}{2}, 2k\pi + \frac{\pi}{2}\right) (k \in \mathbb{Z})$, 对于函数 $f(x) = 7\sin\left(x - \frac{\pi}{6}\right)$, 由 $2k\pi - \frac{\pi}{2} < x - \frac{\pi}{6} < 2k\pi + \frac{\pi}{2}$, 解得 $2k\pi - \frac{\pi}{3} < x < 2k\pi + \frac{2\pi}{3}$. 取 $k=0$, 得一个单调递增区间为 $\left(-\frac{\pi}{3}, \frac{2\pi}{3}\right)$. 因为 $\left(0, \frac{\pi}{2}\right) \subset \left(-\frac{\pi}{3}, \frac{2\pi}{3}\right)$, 故选A.
\end{solution}
\end{question}

\begin{question}[points=5]
已知 $F_1$, $F_2$ 是椭圆 $C: \frac{x^2}{9} + \frac{y^2}{4} = 1$ 的两个焦点, 点 $M$ 在 $C$ 上, 则 $|MF_1| \cdot |MF_2|$ 的最大值为 ( )
\begin{choices}
  \item $13$
  \item $12$
  \item $9$
  \item $6$
\end{choices}
\begin{solution}
\textbf{答案：}$9$

\textbf{分析：}本题通过利用椭圆定义得到 $|MF_1| + |MF_2| = 2a = 6$, 借助基本不等式 $|MF_1| \cdot |MF_2| \leq \left(\frac{|MF_1| + |MF_2|}{2}\right)^2$ 即可得到答案.

\textbf{详解：}由题, $a^2 = 9, b^2 = 4$, 则 $|MF_1| + |MF_2| = 2a = 6$, 所以 $|MF_1| \cdot |MF_2| \leq \left(\frac{|MF_1| + |MF_2|}{2}\right)^2 = 9$ (当且仅当 $|MF_1| = |MF_2| = 3$ 时等号成立), 故选C.
\end{solution}
\end{question}

\begin{question}[points=5]
若 $\tan\theta = -2$, 则 $\frac{\sin\theta(1 + \sin2\theta)}{\sin\theta + \cos\theta} =$ ( )
\begin{choices}
  \item $-\frac{6}{5}$
  \item $-\frac{2}{5}$
  \item $\frac{2}{5}$
  \item $\frac{6}{5}$
\end{choices}
\begin{solution}
\textbf{答案：}$\frac{2}{5}$

\textbf{分析：}将式子进行齐次化处理, 代入 $\tan\theta = -2$ 即可.

\textbf{详解：}$\frac{\sin\theta(1 + \sin2\theta)}{\sin\theta + \cos\theta} = \frac{\sin\theta(\sin^2\theta + \cos^2\theta + 2\sin\theta\cos\theta)}{\sin\theta + \cos\theta} = \sin\theta(\sin\theta + \cos\theta) = \frac{\sin\theta(\sin\theta + \cos\theta)}{\sin^2\theta + \cos^2\theta} = \frac{\tan^2\theta + \tan\theta}{1 + \tan^2\theta} = \frac{4 - 2}{1 + 4} = \frac{2}{5}$, 故选C.
\end{solution}
\end{question}

\begin{question}[points=5]
若过点 $(a,b)$ 可以作曲线 $y = e^x$ 的两条切线, 则 ( )
\begin{choices}
  \item $e^b < a$
  \item $e^a < b$
  \item $0 < a < e^b$
  \item $0 < b < e^a$
\end{choices}
\begin{solution}
\textbf{答案：}$0 < b < e^a$

\textbf{分析：}根据导数几何意义求得切线方程, 再构造函数, 利用导数研究函数图象, 结合图形确定结果.

\textbf{详解：}在曲线 $y = e^x$ 上任取一点 $P(t, e^t)$, 对函数求导得 $y' = e^x$, 曲线在点 $P$ 处的切线方程为 $y - e^t = e^t(x - t)$, 即 $y = e^t x + (1 - t)e^t$. 由题意可知点 $(a,b)$ 在切线上, 得 $b = ae^t + (1 - t)e^t = (a+1-t)e^t$. 令 $f(t) = (a+1-t)e^t$, 则 $f'(t) = (a - t)e^t$. 当 $t < a$ 时 $f'(t) > 0$, $f(t)$ 递增; 当 $t > a$ 时 $f'(t) < 0$, $f(t)$ 递减, 所以 $f(t)_{\max} = f(a) = e^a$. 由题意, 直线 $y = b$ 与曲线 $y = f(t)$ 有两个交点, 则 $b < e^a$. 当 $t < a+1$ 时 $f(t) > 0$, 当 $t > a+1$ 时 $f(t) < 0$, 结合图象知 $0 < b < e^a$, 故选D.
\end{solution}
\end{question}

\begin{question}[points=5]
有6个相同的球, 分别标有数字1,2,3,4,5,6, 从中有放回的随机取两次, 每次取1个球. 甲表示事件“第一次取出的球的数字是1”, 乙表示事件“第二次取出的球的数字是2”, 丙表示事件“两次取出的球的数字之和是8”, 丁表示事件“两次取出的球的数字之和是7”, 则 ( )
\begin{choices}
  \item 甲与丙相互独立
  \item 甲与丁相互独立
  \item 乙与丙相互独立
  \item 丙与丁相互独立
\end{choices}
\begin{solution}
\textbf{答案：}甲与丁相互独立

\textbf{分析：}根据独立事件概率关系逐一判断.

\textbf{详解：}$P(\text{甲}) = \frac{1}{6}$, $P(\text{乙}) = \frac{1}{6}$, $P(\text{丙}) = \frac{5}{36}$, $P(\text{丁}) = \frac{6}{36} = \frac{1}{6}$. $P(\text{甲丙}) = 0 \neq P(\text{甲})P(\text{丙})$, $P(\text{甲丁}) = \frac{1}{36} = P(\text{甲})P(\text{丁})$, $P(\text{乙丙}) = \frac{1}{36} \neq P(\text{乙})P(\text{丙})$, $P(\text{丙丁}) = 0 \neq P(\text{丁})P(\text{丙})$, 故选B.
\end{solution}
\end{question}

\section*{二、多项选择题}
本题共4小题，每小题5分，共20分。在每小题给出的选项中，有多项符合题目要求。全部选对的得5分，部分选对的得2分，有选错的得0分。

\begin{question}[points=5]
有一组样本数据 $x_1, x_2, \ldots, x_n$, 由这组数据得到新样本数据 $y_1, y_2, \ldots, y_n$, 其中 $y_i = x_i + c$ ($i = 1,2,\ldots,n$), $c$ 为非零常数, 则 ( )
\begin{choices}
  \item 两组样本数据的样本平均数相同
  \item 两组样本数据的样本中位数相同
  \item 两组样本数据的样本标准差相同
  \item 两组样本数据的样本极差相同
\end{choices}
\begin{solution}
\textbf{答案：}C, D

\textbf{分析：}A、C利用两组数据的线性关系有 $E(y) = E(x) + c$, $D(y) = D(x)$, 即可判断; 根据中位数、极差的定义, 结合已知线性关系可判断B、D.

\textbf{详解：}A: $E(y) = E(x+c) = E(x) + c$, 故平均数不相同, 错误; B: 若第一组中位数为 $x_i$, 则第二组中位数为 $y_i = x_i + c$, 不相同, 错误; C: $D(y) = D(x) + D(c) = D(x)$, 故方差相同, 正确; D: 极差 $y_{\max} - y_{\min} = (x_{\max} + c) - (x_{\min} + c) = x_{\max} - x_{\min}$, 故极差相同, 正确. 故选CD.
\end{solution}
\end{question}

\begin{question}[points=5]
已知 $O$ 为坐标原点, 点 $P_1(\cos\alpha, \sin\alpha)$, $P_2(\cos\beta, -\sin\beta)$, $P_3(\cos(\alpha+\beta), \sin(\alpha+\beta))$, $A(1,0)$, 则 ( )
\begin{choices}
  \item $|\overrightarrow{OP_1}| = |\overrightarrow{OP_2}|$
  \item $|\overrightarrow{AP_1}| = |\overrightarrow{AP_2}|$
  \item $\overrightarrow{OA} \cdot \overrightarrow{OP_3} = \overrightarrow{OP_1} \cdot \overrightarrow{OP_2}$
  \item $\overrightarrow{OA} \cdot \overrightarrow{OP_1} = \overrightarrow{OP_2} \cdot \overrightarrow{OP_3}$
\end{choices}
\begin{solution}
\textbf{答案：}A, C

\textbf{分析：}A、B写出向量坐标, 利用坐标公式求模; C、D应用向量数量积的坐标表示及两角和差公式化简.

\textbf{详解：}A: $\overrightarrow{OP_1} = (\cos\alpha, \sin\alpha)$, $\overrightarrow{OP_2} = (\cos\beta, -\sin\beta)$, $|\overrightarrow{OP_1}| = \sqrt{\cos^2\alpha + \sin^2\alpha} = 1$, $|\overrightarrow{OP_2}| = 1$, 故正确; B: $|\overrightarrow{AP_1}| = \sqrt{(\cos\alpha-1)^2 + \sin^2\alpha} = \sqrt{2(1-\cos\alpha)} = 2|\sin\frac{\alpha}{2}|$, $|\overrightarrow{AP_2}| = 2|\sin\frac{\beta}{2}|$, 不一定相等, 错误; C: $\overrightarrow{OA} \cdot \overrightarrow{OP_3} = \cos(\alpha+\beta)$, $\overrightarrow{OP_1} \cdot \overrightarrow{OP_2} = \cos\alpha\cos\beta + \sin\alpha(-\sin\beta) = \cos(\alpha+\beta)$, 正确; D: $\overrightarrow{OA} \cdot \overrightarrow{OP_1} = \cos\alpha$, $\overrightarrow{OP_2} \cdot \overrightarrow{OP_3} = \cos\beta\cos(\alpha+\beta) + (-\sin\beta)\sin(\alpha+\beta) = \cos(\alpha+2\beta)$, 不一定相等, 错误. 故选AC.
\end{solution}
\end{question}

\begin{question}[points=5]
已知点 $P$ 在圆 $(x-5)^2 + (y-5)^2 = 16$ 上, 点 $A(4,0)$, $B(0,2)$, 则 ( )
\begin{choices}
  \item 点 $P$ 到直线 $AB$ 的距离小于 $10$
  \item 点 $P$ 到直线 $AB$ 的距离大于 $2$
  \item 当 $\angle PBA$ 最小时, $|PB| = 3\sqrt{2}$
  \item 当 $\angle PBA$ 最大时, $|PB| = 3\sqrt{2}$
\end{choices}
\begin{solution}
\textbf{答案：}A, C, D

\textbf{分析：}计算出圆心到直线 $AB$ 的距离, 得出点 $P$ 到直线 $AB$ 的距离的取值范围; 当 $\angle PBA$ 最大或最小时, $PB$ 与圆 $M$ 相切, 利用勾股定理判断.

\textbf{详解：}圆 $(x-5)^2+(y-5)^2=16$ 的圆心 $M(5,5)$, 半径 $4$, 直线 $AB$ 方程: $\frac{x}{4}+\frac{y}{2}=1$, 即 $x+2y-4=0$, 圆心到直线距离 $d = \frac{|5+10-4|}{\sqrt{1^2+2^2}} = \frac{11}{\sqrt{5}} = \frac{11\sqrt{5}}{5} > 4$, 故点 $P$ 到直线距离最小值 $\frac{11\sqrt{5}}{5} - 4 < 2$, 最大值 $\frac{11\sqrt{5}}{5} + 4 < 10$, 故A正确, B错误. 当 $\angle PBA$ 最大或最小时, $PB$ 与圆 $M$ 相切, $|BM| = \sqrt{(0-5)^2+(2-5)^2} = \sqrt{34}$, $|MP| = 4$, 由勾股定理 $|BP| = \sqrt{34-16} = 3\sqrt{2}$, 故CD正确. 故选ACD.
\end{solution}
\end{question}

\begin{question}[points=5]
在正三棱柱 $ABC-A_1B_1C_1$ 中, $AB = AA_1 = 1$, 点 $P$ 满足 $\overrightarrow{BP} = \lambda \overrightarrow{BC} + \mu \overrightarrow{BB_1}$, 其中 $\lambda \in [0,1]$, $\mu \in [0,1]$, 则 ( )
\begin{choices}
  \item 当 $\lambda = 1$ 时, $\triangle AB_1P$ 的周长为定值
  \item 当 $\mu = 1$ 时, 三棱锥 $P-A_1BC$ 的体积为定值
  \item 当 $\lambda = \frac{1}{2}$ 时, 有且仅有一个点 $P$, 使得 $A_1P \perp BP$
  \item 当 $\mu = \frac{1}{2}$ 时, 有且仅有一个点 $P$, 使得 $A_1B \perp$ 平面 $AB_1P$
\end{choices}
\begin{solution}
\textbf{答案：}B, D

\textbf{分析：}将向量关系转化为点的轨迹, 通过建系或几何性质判断.

\textbf{详解：}A: $\lambda = 1$ 时 $P$ 在 $CC_1$ 上, $\triangle AB_1P$ 周长不是定值, 错误; B: $\mu = 1$ 时 $P$ 在 $B_1C_1$ 上, $BC \parallel B_1C_1$, 故 $P$ 到平面 $A_1BC$ 距离为定值, 体积为定值, 正确; C: $\lambda = \frac{1}{2}$ 时 $P$ 在 $BC$ 中点与 $B_1C_1$ 中点的连线上, 建系得 $A_1P \perp BP$ 有 $\mu=0$ 或 $1$, 两个点, 错误; D: $\mu = \frac{1}{2}$ 时 $P$ 在 $BB_1$ 中点与 $CC_1$ 中点的连线上, 建系得 $A_1B \perp$ 平面 $AB_1P$ 时 $P$ 与 $C_1$ 重合, 唯一, 正确. 故选BD.
\end{solution}
\end{question}

\section*{三、填空题}
本题共4小题，每小题5分，共20分。

\begin{question}[points=5]
已知函数 $f(x) = x^3(a \cdot 2^x - 2^{-x})$ 是偶函数, 则 $a =$ .
\fillin[{$1$}]
\begin{solution}
\textbf{答案：}$1$

\textbf{分析：}利用偶函数的定义可求参数 $a$ 的值.

\textbf{详解：}因为 $f(x) = x^3(a \cdot 2^x - 2^{-x})$, 故 $f(-x) = -x^3(a \cdot 2^{-x} - 2^x)$. 由 $f(-x) = f(x)$ 得 $x^3(a \cdot 2^x - 2^{-x}) = -x^3(a \cdot 2^{-x} - 2^x)$, 整理得 $(a-1)(2^x + 2^{-x}) = 0$, 所以 $a = 1$.
\end{solution}
\end{question}

\begin{question}[points=5]
已知 $O$ 为坐标原点, 抛物线 $C: y^2 = 2px \ (p > 0)$ 的焦点为 $F$, $P$ 为 $C$ 上一点, $PF$ 与 $x$ 轴垂直, $Q$ 为 $x$ 轴上一点, 且 $PQ \perp OP$, 若 $|FQ| = 6$, 则 $C$ 的准线方程为.
\fillin[{$x = -\frac{3}{2}$}]
\begin{solution}
\textbf{答案：}$x = -\frac{3}{2}$

\textbf{分析：}先用坐标表示 $P$, $Q$, 再根据向量垂直坐标关系列方程, 解得 $p$.

\textbf{详解：}不妨设 $P\left(\frac{p}{2}, p\right)$, 则 $Q\left(6 + \frac{p}{2}, 0\right)$, $\overrightarrow{PQ} = (6, -p)$. 因为 $PQ \perp OP$, 所以 $\frac{p}{2} \times 6 - p^2 = 0$, 解得 $p = 3$ (舍去 $p=0$), 故准线方程为 $x = -\frac{p}{2} = -\frac{3}{2}$.
\end{solution}
\end{question}

\begin{question}[points=5]
函数 $f(x) = |2x - 1| - 2\ln x$ 的最小值为.
\fillin[{$1$}]
\begin{solution}
\textbf{答案：}$1$

\textbf{分析：}由解析式知定义域为 $(0, +\infty)$, 分段讨论并利用导数研究单调性, 求最小值.

\textbf{详解：}定义域 $(0, +\infty)$. 当 $0 < x \leq \frac{1}{2}$ 时, $f(x) = 1 - 2x - 2\ln x$, 单调递减; 当 $\frac{1}{2} < x \leq 1$ 时, $f(x) = 2x - 1 - 2\ln x$, $f'(x) = 2 - \frac{2}{x} \leq 0$, 递减; 当 $x > 1$ 时, $f(x) = 2x - 1 - 2\ln x$, $f'(x) = 2 - \frac{2}{x} > 0$, 递增. 故 $f(x)$ 在 $(0,1]$ 递减, 在 $[1, +\infty)$ 递增, 最小值为 $f(1) = 1$.
\end{solution}
\end{question}

\begin{question}[points=5]
某校学生在研究民间剪纸艺术时, 发现剪纸时经常会沿纸的某条对称轴把纸对折. 规格为 $20\,\text{dm} \times 12\,\text{dm}$ 的长方形纸, 对折1次共可以得到 $10\,\text{dm} \times 12\,\text{dm}$, $20\,\text{dm} \times 6\,\text{dm}$ 两种规格的图形, 它们的面积之和 $S_1 = 240\,\text{dm}^2$; 对折2次共可以得到 $5\,\text{dm} \times 12\,\text{dm}$, $10\,\text{dm} \times 6\,\text{dm}$, $20\,\text{dm} \times 3\,\text{dm}$ 三种规格的图形, 它们的面积之和 $S_2 = 180\,\text{dm}^2$; 以此类推, 则对折4次共可以得到不同规格图形的种数为; 如果对折 $n$ 次, 那么 $\sum_{k=1}^{n} S_k =$ $\text{dm}^2$.
\fillin[{$5$ ; $720 - \frac{15(n+3)}{2^{n-4}}$}]
\begin{solution}
\textbf{答案：}$5$ ; $720 - \frac{15(n+3)}{2^{n-4}}$

\textbf{分析：}列举对折4次可得规格种类数; 找出 $S_n$ 的表达式, 再利用错位相减法求和.

\textbf{详解：}对折4次可得到: $\frac{5}{4}\times12$, $\frac{5}{2}\times6$, $5\times3$, $10\times\frac{3}{2}$, $20\times\frac{3}{4}$, 共5种. $S_1=2\times120$, $S_2=3\times60$, $S_3=4\times30$, $S_4=5\times15$, 一般地 $S_n = \frac{120(n+1)}{2^{n-1}}$. 设 $S = \sum_{k=1}^n S_k = \sum_{k=1}^n \frac{120(k+1)}{2^{k-1}}$. 用错位相减法得 $S = 720 - \frac{15(n+3)}{2^{n-4}}$.
\end{solution}
\end{question}

\section*{四、解答题}
本题共6小题，共70分。解答应写出文字说明、证明过程或演算步骤。

\begin{question}[points=10]
已知数列 $\{a_n\}$ 满足 $a_1 = 1$, $a_{n+1} = \begin{cases} a_n + 1, & n \text{为奇数}, \\ a_n + 2, & n \text{为偶数}. \end{cases}$

(1) 记 $b_n = a_{2n}$, 写出 $b_1$, $b_2$, 并求数列 $\{b_n\}$ 的通项公式;

(2) 求 $\{a_n\}$ 的前20项和.
\begin{solution}
\textbf{答案：}(1) $b_1 = 2$, $b_2 = 5$, $b_n = 3n - 1$; (2) $300$

\textbf{分析：}(1) 根据递推关系可得 $b_{n+1} = b_n + 3$, 从而可求通项. (2) 利用奇数项与偶数项的关系求和.

\textbf{详解：}(1) $b_1 = a_2 = a_1 + 1 = 2$, $b_2 = a_4 = a_3 + 1 = (a_2 + 2) + 1 = 5$. 由递推得 $a_{2k+2} = a_{2k+1} + 1$, $a_{2k+1} = a_{2k} + 2$, 故 $a_{2k+2} = a_{2k} + 3$, 即 $b_{n+1} = b_n + 3$, 所以 $\{b_n\}$ 是等差数列, $b_n = 2 + (n-1)\times 3 = 3n - 1$.
(2) $S_{20} = a_1 + a_2 + \cdots + a_{20}$. 因为 $a_1 = a_2 - 1$, $a_3 = a_4 - 1$, $\ldots$, $a_{19} = a_{20} - 1$, 所以 $S_{20} = 2(a_2 + a_4 + \cdots + a_{20}) - 10 = 2(b_1 + b_2 + \cdots + b_{10}) - 10 = 2\left(10\times 2 + \frac{9\times 10}{2}\times 3\right) - 10 = 300$.
\end{solution}
\end{question}

\begin{question}[points=12]
某学校组织“一带一路”知识竞赛, 有A, B两类问题. 每位参加比赛的同学先在两类问题中选择一类并从中随机抽取一个问题回答, 若回答错误则该同学比赛结束; 若回答正确则从另一类问题中再随机抽取一个问题回答, 无论回答正确与否, 该同学比赛结束. A类问题中的每个问题回答正确得20分, 否则得0分; B类问题中的每个问题回答正确得80分, 否则得0分. 已知小明能正确回答A类问题的概率为0.8, 能正确回答B类问题的概率为0.6, 且能正确回答问题的概率与回答次序无关.

(1) 若小明先回答A类问题, 记 $X$ 为小明的累计得分, 求 $X$ 的分布列;

(2) 为使累计得分的期望最大, 小明应选择先回答哪类问题? 并说明理由.
\begin{solution}
\textbf{答案：}(1) 分布列见解析; (2) 应先回答B类问题, 理由见解析.

\textbf{分析：}(1) 分析 $X$ 的所有可能取值, 分别求概率; (2) 类似求出先回答B类问题的期望, 比较得出.

\textbf{详解：}(1) $X$ 的可能取值为0, 20, 100. $P(X=0) = 1-0.8 = 0.2$; $P(X=20) = 0.8 \times (1-0.6) = 0.32$; $P(X=100) = 0.8 \times 0.6 = 0.48$. 分布列为:
| $X$ | 0 | 20 | 100 |
|-----|---|---|-----|
| $P$ | 0.2 | 0.32 | 0.48 |
$E(X) = 0\times0.2 + 20\times0.32 + 100\times0.48 = 54.4$.
(2) 若先回答B类, 设 $Y$ 为累计得分, $Y$ 的可能取值为0, 80, 100. $P(Y=0) = 1-0.6 = 0.4$; $P(Y=80) = 0.6\times(1-0.8) = 0.12$; $P(Y=100) = 0.6\times0.8 = 0.48$. $E(Y) = 0\times0.4 + 80\times0.12 + 100\times0.48 = 57.6$. 因为 $54.4 < 57.6$, 所以小明应先回答B类问题.
\end{solution}
\end{question}

\begin{question}[points=12]
记 $\triangle ABC$ 的内角 $A$, $B$, $C$ 的对边分别为 $a$, $b$, $c$. 已知 $b^2 = ac$, 点 $D$ 在边 $AC$ 上, $BD \sin\angle ABC = a \sin C$.

(1) 证明: $BD = b$;

(2) 若 $AD = 2DC$, 求 $\cos\angle ABC$.
\begin{solution}
\textbf{答案：}(1) 证明见解析; (2) $\cos\angle ABC = \frac{7}{12}$

\textbf{分析：}(1) 利用正弦定理的边角关系得到 $BD = \frac{ac}{b}$, 结合 $b^2 = ac$ 得证. (2) 由题设及余弦定理, 利用 $\angle ADB + \angle CDB = \pi$ 建立方程, 解出 $a,b,c$ 关系, 再求余弦值.

\textbf{详解：}(1) 由正弦定理, $\frac{a}{\sin\angle ABC} = \frac{c}{\sin C}$, 故 $\frac{\sin C}{\sin\angle ABC} = \frac{c}{a}$. 又 $BD = \frac{a\sin C}{\sin\angle ABC}$, 所以 $BD = \frac{ac}{b}$. 因为 $b^2 = ac$, 所以 $BD = b$.
(2) 设 $BD = b$, $AD = \frac{2b}{3}$, $DC = \frac{b}{3}$. 在 $\triangle ABD$ 和 $\triangle CBD$ 中分别应用余弦定理, 并利用 $\cos\angle ADB = -\cos\angle CDB$, 可得 $2a^2 + c^2 = \frac{11b^2}{3}$. 结合 $b^2 = ac$, 解得 $\frac{a^2}{b^2} = \frac{1}{3}$ 或 $\frac{3}{2}$. 当 $\frac{a^2}{b^2} = \frac{1}{3}$ 时, $\cos\angle ABC = \frac{a^2 + c^2 - b^2}{2ac} = \frac{7}{6} > 1$, 舍去; 当 $\frac{a^2}{b^2} = \frac{3}{2}$ 时, $\cos\angle ABC = \frac{7}{12}$, 符合题意.
\end{solution}
\end{question}

\begin{question}[points=12]
如图, 在三棱锥 $A-BCD$ 中, 平面 $ABD \perp$ 平面 $BCD$, $AB = AD$, $O$ 为 $BD$ 的中点.

(1) 证明: $OA \perp CD$;

(2) 若 $\triangle OCD$ 是边长为1的等边三角形, 点 $E$ 在棱 $AD$ 上, $DE = 2EA$, 且二面角 $E-BC-D$ 的大小为 $45^\circ$, 求三棱锥 $A-BCD$ 的体积.

\begin{center}
  \includegraphics[width=.38\linewidth]{question-20-1.jpg}
\end{center}
\begin{solution}
\textbf{答案：}(1) 证明见解析; (2) $\frac{\sqrt{3}}{6}$

\textbf{分析：}(1) 利用面面垂直性质定理得 $AO \perp$ 平面 $BCD$, 从而 $AO \perp CD$. (2) 作出二面角平面角, 求得高, 再求体积.

\textbf{详解：}(1) 因为 $AB = AD$, $O$ 为 $BD$ 中点, 所以 $AO \perp BD$. 又平面 $ABD \perp$ 平面 $BCD$, 交线为 $BD$, $AO \subset$ 平面 $ABD$, 所以 $AO \perp$ 平面 $BCD$. 因为 $CD \subset$ 平面 $BCD$, 所以 $AO \perp CD$.
(2) 作 $EF \perp BD$ 于 $F$, 作 $FM \perp BC$ 于 $M$, 连接 $EM$. 由(1)知 $AO \perp$ 平面 $BCD$, 故 $EF \perp$ 平面 $BCD$, 则 $EF \perp BC$. 又 $FM \perp BC$, 所以 $BC \perp$ 平面 $EFM$, 故 $\angle EMF$ 为二面角 $E-BC-D$ 的平面角, $\angle EMF = 45^\circ$. 由 $\triangle OCD$ 为正三角形, $O$ 为 $BD$ 中点, 可得 $BD = 2$, $BC = \sqrt{3}$. 由 $DE = 2EA$ 得 $DF = \frac{2}{3}DO = \frac{2}{3}$, $BF = \frac{4}{3}$. $FM = \frac{1}{2}BF = \frac{2}{3}$, 在等腰直角三角形 $EFM$ 中, $EF = FM = \frac{2}{3}$, 故 $AO = \frac{2}{3} \times \frac{3}{2} = 1$? 实际上 $AO = 1$ (由相似或比例得到). 所以 $V_{A-BCD} = \frac{1}{3} \times AO \times S_{\triangle BCD} = \frac{1}{3} \times 1 \times \frac{1}{2} \times 2 \times \sqrt{3} = \frac{\sqrt{3}}{6}$.
\end{solution}
\end{question}

\begin{question}[points=12]
在平面直角坐标系 $xOy$ 中, 已知点 $F_1(-\sqrt{17}, 0)$, $F_2(\sqrt{17}, 0)$, $|MF_1| - |MF_2| = 2$, 点 $M$ 的轨迹为 $C$.

(1) 求 $C$ 的方程;

(2) 设点 $T$ 在直线 $x = \frac{1}{2}$ 上, 过 $T$ 的两条直线分别交 $C$ 于 $A$, $B$ 两点和 $P$, $Q$ 两点, 且 $|TA| \cdot |TB| = |TP| \cdot |TQ|$, 求直线 $AB$ 的斜率与直线 $PQ$ 的斜率之和.
\begin{solution}
\textbf{答案：}(1) $x^2 - \frac{y^2}{16} = 1 \ (x \geq 1)$; (2) $0$

\textbf{分析：}(1) 利用双曲线定义求得轨迹方程. (2) 设直线方程, 与双曲线联立, 利用弦长公式及条件建立斜率关系.

\textbf{详解：}(1) 由 $|MF_1| - |MF_2| = 2 < |F_1F_2| = 2\sqrt{17}$, 得轨迹为双曲线右支. $2a = 2$, $a = 1$, $c = \sqrt{17}$, $b^2 = c^2 - a^2 = 16$, 方程 $x^2 - \frac{y^2}{16} = 1 \ (x \geq 1)$.
(2) 设 $T(\frac{1}{2}, t)$, 直线 $AB$ 斜率 $k_1$, 方程 $y - t = k_1(x - \frac{1}{2})$, 与双曲线联立, 利用韦达定理得 $|TA| \cdot |TB| = \frac{(t^2+12)(1+k_1^2)}{k_1^2-16}$. 同理, 直线 $PQ$ 斜率 $k_2$, $|TP| \cdot |TQ| = \frac{(t^2+12)(1+k_2^2)}{k_2^2-16}$. 由 $|TA| \cdot |TB| = |TP| \cdot |TQ|$ 得 $\frac{1+k_1^2}{k_1^2-16} = \frac{1+k_2^2}{k_2^2-16}$, 整理得 $(k_1^2 - k_2^2)(1+16)=0$, 即 $k_1^2 = k_2^2$. 由于 $k_1 \neq k_2$ (否则为同一直线), 故 $k_1 + k_2 = 0$.
\end{solution}
\end{question}

\begin{question}[points=12]
已知函数 $f(x) = x(1 - \ln x)$.

(1) 讨论 $f(x)$ 的单调性;

(2) 设 $a$, $b$ 为两个不相等的正数, 且 $b\ln a - a\ln b = a - b$, 证明: $2 < \frac{1}{a} + \frac{1}{b} < e$.
\begin{solution}
\textbf{答案：}(1) $f(x)$ 的递增区间为 $(0,1)$, 递减区间为 $(1,+\infty)$; (2) 证明见解析.

\textbf{分析：}(1) 求导判断符号. (2) 变形为 $\frac{1}{a}(1-\ln\frac{1}{a}) = \frac{1}{b}(1-\ln\frac{1}{b})$, 设 $x_1 = \frac{1}{a}$, $x_2 = \frac{1}{b}$, 则 $f(x_1) = f(x_2)$, 利用极值点偏移证明 $x_1 + x_2 > 2$, 再构造函数证明 $x_1 + x_2 < e$.

\textbf{详解：}(1) $f'(x) = -\ln x$, 当 $x \in (0,1)$ 时 $f'(x) > 0$, 当 $x \in (1,+\infty)$ 时 $f'(x) < 0$, 故 $f(x)$ 在 $(0,1)$ 递增, 在 $(1,+\infty)$ 递减.
(2) 由 $b\ln a - a\ln b = a - b$ 得 $\frac{\ln a + 1}{a} = \frac{\ln b + 1}{b}$, 即 $\frac{1}{a}(1-\ln\frac{1}{a}) = \frac{1}{b}(1-\ln\frac{1}{b})$. 令 $x_1 = \frac{1}{a}$, $x_2 = \frac{1}{b}$, 则 $f(x_1) = f(x_2)$, 不妨设 $0 < x_1 < 1 < x_2$. 先证 $x_1 + x_2 > 2$: 若 $x_2 \geq 2$, 显然成立; 若 $1 < x_2 < 2$, 设 $g(x) = f(x) - f(2-x)$ $(1 < x < 2)$, $g'(x) = -\ln[x(2-x)] > 0$, $g(x)$ 递增, $g(x) > g(1) = 0$, 故 $f(x_2) > f(2-x_2)$, 即 $f(x_1) > f(2-x_2)$, 由单调性得 $x_1 < 2 - x_2$, 即 $x_1 + x_2 < 2$? 需要印证: 实际上 $f(x_1)=f(x_2)$, 由 $f(x_2) > f(2-x_2)$ 得 $f(x_1) > f(2-x_2)$, 因为 $0<x_1<1$, $0<2-x_2<1$, 且 $f$ 在 $(0,1)$ 递增, 所以 $x_1 > 2 - x_2$, 即 $x_1 + x_2 > 2$. 再证 $x_1 + x_2 < e$: 设 $x_2 = t x_1 (t>1)$, 由 $f(x_1)=f(x_2)$ 得 $\ln x_1 = \frac{t-1 - t\ln t}{t-1}$. 要证 $x_1 + x_2 < e$, 即 $\ln x_1 + \ln(t+1) < 1$, 代入得 $\frac{t-1 - t\ln t}{t-1} + \ln(t+1) < 1$, 即 $(t-1)\ln(t+1) - t\ln t < 0$. 令 $S(t) = (t-1)\ln(t+1) - t\ln t$, $S'(t) = \ln(1+\frac{1}{t}) - \frac{2}{t+1}$. 利用不等式 $\ln(1+x) \leq x$ 得 $\ln(1+\frac{1}{t}) < \frac{1}{t} < \frac{2}{t+1}$ (当 $t>1$), 所以 $S'(t) < 0$, $S(t)$ 递减, $S(t) < S(1) = 0$, 故 $x_1 + x_2 < e$ 成立. 综上, $2 < \frac{1}{a} + \frac{1}{b} < e$.
\end{solution}
\end{question}

